%=========================
% Sections/05-KhaiThacYeuCau.tex
% Chương 5 – Khai thác yêu cầu
%=========================
\section{Khai thác yêu cầu}

Chương này xác định phạm vi và các yêu cầu của website thương mại điện tử tích hợp trợ lý mua sắm đa phương thức. Trọng tâm là tìm kiếm sản phẩm bằng ngôn ngữ tự nhiên, hình ảnh và hội thoại nhiều lượt; các chức năng thương mại điện tử cung cấp quy trình xem và mua sản phẩm. Yêu cầu là cơ sở cho thiết kế kiến trúc, cơ sở dữ liệu, pipeline AI và kiểm thử.

\subsection{Tác nhân hệ thống}

\begin{longtable}{|p{0.20\linewidth}|p{0.34\linewidth}|p{0.34\linewidth}|}
\caption{Các tác nhân của hệ thống}\label{tab:system-actors}\\
\hline
\textbf{Tác nhân} & \textbf{Mục tiêu} & \textbf{Tác vụ chính} \\
\hline
\endfirsthead
\hline
\textbf{Tác nhân} & \textbf{Mục tiêu} & \textbf{Tác vụ chính} \\
\hline
\endhead
Khách hàng & Tìm kiếm, đánh giá và mua sản phẩm. & Tài khoản; tìm kiếm text/image; hội thoại; lọc và xem sản phẩm; giỏ hàng; đặt hàng COD; theo dõi đơn hàng; đánh giá. \\
\hline
Quản trị viên/Nhân viên & Vận hành dữ liệu và nghiệp vụ thương mại điện tử. & Quản lý sản phẩm, hình ảnh, tồn kho, đơn hàng, tài khoản, nội dung và phân quyền. \\
\hline
Hệ thống AI & Xử lý truy vấn và sinh phản hồi có ngữ cảnh. & Phân tích ý định; embedding; Qdrant retrieval; lọc metadata; fusion/rerank; bộ nhớ hội thoại; log; đồng bộ vector. \\
\hline
\end{longtable}

Re-index và theo dõi retrieval là tác vụ vận hành nội bộ, không phải actor người dùng độc lập.

\subsection{Yêu cầu chức năng}

\subsubsection{Trợ lý mua sắm đa phương thức}

\begin{longtable}{|p{0.12\linewidth}|p{0.25\linewidth}|p{0.45\linewidth}|p{0.10\linewidth}|}
\caption{Yêu cầu chức năng AI}\label{tab:ai-functional-requirements}\\
\hline
\textbf{Mã} & \textbf{Chức năng} & \textbf{Mô tả} & \textbf{Ưu tiên} \\
\hline
\endfirsthead
\hline
\textbf{Mã} & \textbf{Chức năng} & \textbf{Mô tả} & \textbf{Ưu tiên} \\
\hline
\endhead
FR-AI-01 & Tìm kiếm text & Tạo text embedding và truy hồi sản phẩm phù hợp theo ngữ nghĩa. & Cao \\
\hline
FR-AI-02 & Tìm kiếm image & Nhận ảnh, tạo image embedding và tìm sản phẩm tương tự. & Cao \\
\hline
FR-AI-03 & Hybrid search & Kết hợp ảnh tham chiếu với điều kiện văn bản. & Cao \\
\hline
FR-AI-04 & Phân tích truy vấn & Trích xuất loại sản phẩm, màu, kích thước, thương hiệu và khoảng giá. & Cao \\
\hline
FR-AI-05 & Lọc metadata & Lọc theo giá, danh mục, màu sắc, thương hiệu và tồn kho. & Cao \\
\hline
FR-AI-06 & Hội thoại nhiều lượt & Duy trì ngữ cảnh, ràng buộc và lựa chọn trước đó. & Cao \\
\hline
FR-AI-07 & Phản hồi có căn cứ & Chỉ sinh phản hồi từ dữ liệu sản phẩm được truy hồi. & Cao \\
\hline
FR-AI-08 & Fusion và rerank & Hợp nhất các nhánh text/image và sắp xếp ứng viên. & Cao \\
\hline
FR-AI-09 & Gợi ý tương tự & Đề xuất sản phẩm liên quan dựa trên độ tương đồng. & TB \\
\hline
FR-AI-10 & Đồng bộ chỉ mục & Cập nhật hoặc xóa embedding/payload Qdrant khi sản phẩm thay đổi. & Cao \\
\hline
FR-AI-11 & Lịch sử và log & Lưu checkpoint, trạng thái phiên, latency và log cần thiết. & TB \\
\hline
FR-AI-12 & Phản hồi dự phòng & thương mại điện tử cơ bản vẫn hoạt động khi dịch vụ AI tạm lỗi. & Cao \\
\hline
\end{longtable}

YOLO chỉ là bước crop tùy chọn ở môi trường phát triển; tìm kiếm đa ngôn ngữ và mô hình local là hướng mở rộng.

\subsubsection{Website thương mại điện tử}

\begin{longtable}{|p{0.14\linewidth}|p{0.26\linewidth}|p{0.42\linewidth}|p{0.10\linewidth}|}
\caption{Yêu cầu chức năng thương mại điện tử}\label{tab:ecommerce-functional-requirements}\\
\hline
\textbf{Mã} & \textbf{Chức năng} & \textbf{Mô tả} & \textbf{Ưu tiên} \\
\hline
\endfirsthead
\hline
\textbf{Mã} & \textbf{Chức năng} & \textbf{Mô tả} & \textbf{Ưu tiên} \\
\hline
\endhead
EC-FR-01 & Tài khoản & Đăng ký, đăng nhập, cập nhật thông tin và đổi mật khẩu. & Cao \\
\hline
EC-FR-02 & Sản phẩm/giỏ hàng & Xem, lọc sản phẩm; thêm, sửa, xóa sản phẩm trong giỏ. & Cao \\
\hline
EC-FR-03 & Đặt hàng COD & Xác nhận nhận hàng, tạo đơn và thanh toán khi nhận hàng. & Cao \\
\hline
EC-FR-04 & Đơn hàng/đánh giá & Theo dõi đơn hàng và đánh giá sản phẩm đã mua. & TB \\
\hline
EC-FR-05 & Quản trị & Quản lý sản phẩm, ảnh, tồn kho, đơn hàng, người dùng, blog và dashboard. & Cao \\
\hline
\end{longtable}

\subsection{Yêu cầu dữ liệu}

\begin{longtable}{|p{0.22\linewidth}|p{0.52\linewidth}|p{0.16\linewidth}|}
\caption{Các nhóm dữ liệu cần quản lý}\label{tab:data-requirements}\\
\hline
\textbf{Nhóm} & \textbf{Nội dung và ràng buộc} & \textbf{Lưu trữ} \\
\hline
\endfirsthead
\hline
\textbf{Nhóm} & \textbf{Nội dung và ràng buộc} & \textbf{Lưu trữ} \\
\hline
\endhead
Sản phẩm & Tên, mô tả, danh mục, màu, kích thước, giá, tồn kho và URL ảnh. & MySQL \\
\hline
Hình ảnh & Kiểm tra định dạng và kích thước trước khi embedding. & Object storage/URL \\
\hline
Embedding/payload & Text/image vector, mã sản phẩm và metadata cho search/filter/rerank. & Qdrant \\
\hline
Hội thoại & Tin nhắn, tóm tắt, ràng buộc và checkpoint theo phiên. & Redis \\
\hline
Giao dịch/log & Tài khoản, giỏ hàng, đơn COD, đánh giá và log vận hành. & MySQL/monitoring \\
\hline
\end{longtable}

MySQL là nguồn dữ liệu nghiệp vụ chính. Khi sản phẩm hoặc hình ảnh thay đổi, vector Qdrant phải được cập nhật hoặc xóa; lỗi đồng bộ không được làm mất dữ liệu giao dịch.

\subsection{Yêu cầu phi chức năng}

\begin{longtable}{|p{0.14\linewidth}|p{0.20\linewidth}|p{0.50\linewidth}|p{0.08\linewidth}|}
\caption{Yêu cầu phi chức năng}\label{tab:non-functional-requirements}\\
\hline
\textbf{Mã} & \textbf{Nhóm} & \textbf{Mục tiêu và ràng buộc} & \textbf{ƯT} \\
\hline
\endfirsthead
\hline
\textbf{Mã} & \textbf{Nhóm} & \textbf{Mục tiêu và ràng buộc} & \textbf{ƯT} \\
\hline
\endhead
NFR-01 & Hiệu năng & Với 500--2.000 sản phẩm và dưới 50 người dùng đồng thời, đo theo P95: API thương mại điện tử không quá 3 giây; Text Search, Image Search, Hybrid Search và Assistant không quá 5 giây. Tỷ lệ lỗi của bài tải không quá 1\%. & Cao \\
\hline
NFR-02 & Bảo mật & Phân quyền theo vai trò; mật khẩu băm bằng BCrypt; kiểm tra quyền API quản trị và tệp ảnh. & Cao \\
\hline
NFR-03 & Khả dụng & Xem sản phẩm, giỏ hàng và đơn hàng vẫn dùng được khi AI tạm gián đoạn. & Cao \\
\hline
NFR-04 & Độ tin cậy & Dữ liệu giao dịch ghi chính xác trong MySQL; lỗi vector không làm mất dữ liệu nghiệp vụ. & Cao \\
\hline
NFR-05 & Bảo trì/mở rộng & Mô-đun hóa Frontend, Spring Boot, AI Service và adapter để thay thế mô hình hoặc nhà cung cấp. & TB \\
\hline
NFR-06 & Tương thích & Hỗ trợ Chrome, Edge, Firefox, laptop và điện thoại; Backend trao đổi JSON. & TB \\
\hline
NFR-07 & Quan sát & Ghi lỗi, latency, retrieval, rerank và checkpoint đủ để chẩn đoán, không lộ dữ liệu nhạy cảm. & Cao \\
\hline
NFR-08 & Nhất quán & MySQL là nguồn nghiệp vụ; Qdrant đồng bộ sản phẩm; Redis lưu cache/checkpoint, không thay thế giao dịch. & Cao \\
\hline
NFR-09 & Tài nguyên/AI & Production ưu tiên dịch vụ managed; YOLO và model local chỉ bật khi cần. Đo latency và top-k relevance trên truy vấn text/image/hybrid. & TB \\
\hline
\end{longtable}

\subsection*{Kết chương}

Chương này đã xác định tác nhân, chức năng, dữ liệu và ràng buộc phi chức năng cho hệ thống. Phạm vi trọng tâm là text search, image search, hybrid retrieval, lọc metadata và hội thoại có ngữ cảnh; các chức năng thương mại điện tử cung cấp nền tảng để xem và mua sản phẩm. Đây là cơ sở cho các chương thiết kế và kiểm thử tiếp theo.
